% ---- pg-018 ----
\[
\begin{bmatrix} x_{11} \\ x_{12} \end{bmatrix},\quad
\begin{bmatrix} x_{21} \\ x_{22} \end{bmatrix},\quad \cdots,\quad
\begin{bmatrix} x_{n1} \\ x_{n2} \end{bmatrix}
\]

其中 $x_{j1}$ 和 $x_{j2}$ 是由第 $j$ 次试验项观测的($j=1,2,\cdots,n$).变量 1 和 2 的测量值的线性结合由样本协方差给出
\[
s_{12} = \frac{1}{n}\sum_{j=1}^{n}(x_{j1}-\bar{x}_1)(x_{j2}-\bar{x}_2)
\]

或是它们各自的平均值的偏差之积的平均值.若对一个变量的大观测值与对另一变量的大观测值一起出现,而小值也一起出现,$s_{12}$ 将为正数.若一个变量的大值与另一变量的小值一起出现,$s_{12}$ 将为负数.若两个变量的值间没有什么特别的联系,$s_{12}$ 将近似等于零.

\mvsub{样本协方差}
\phantomsection\label{eq:1-4}
\begin{equation}
s_{ik} = \frac{1}{n}\sum_{j=1}^{n}(x_{ji}-\bar{x}_i)(x_{jk}-\bar{x}_k) \qquad i=1,2,\cdots,p, k=1,2,\cdots,p \tag{1-4}
\end{equation}

度量第 $i$ 个和第 $k$ 个变量间的结合.我们注意到,当 $i=k$ 时,协方差就简化为样本方差.此外,对所有的 $i$ 和 $k$ 都有 $s_{ik}=s_{ki}$.

在此,描述统计量最后一个要考虑的是样本相关系数(或皮尔逊积矩相关系数;参见文献[14]).两个变量间线性结合的这个度量不依赖于测量单位,第 $i$ 个和第 $k$ 个变量的样本相关系数定义为
\phantomsection\label{eq:1-5}
\begin{equation}
r_{ik} = \frac{s_{ik}}{\sqrt{s_{ii}}\sqrt{s_{kk}}} = \frac{\sum\limits_{j=1}^{n}(x_{ji}-\bar{x}_i)(x_{jk}-\bar{x}_k)}{\sqrt{\sum\limits_{j=1}^{n}(x_{ji}-\bar{x}_i)^2}\sqrt{\sum\limits_{j=1}^{n}(x_{jk}-\bar{x}_k)^2}} \tag{1-5}
\end{equation}

其中 $i=1,2,\cdots,p$ 和 $k=1,2,\cdots,p$.注意对所有 $i,k$ 都有 $r_{ik}=r_{ki}$.

样本相关系数是样本协方差的一个标准化形式,其中样本方差的平方根的乘积提供了标准化.请注意,对于 $s_{ii},s_{kk}$ 和 $s_{ik}$ 不论是用 $n$ 或 $n-1$ 作为除数,$r_{ik}$ 的值都是相同的.

样本相关系数 $r_{ik}$ 也可看做是一个样本协方差.假设初始值 $x_{ji},x_{jk}$ 由标准化值 $(x_{ji}-\bar{x}_i)/\sqrt{s_{ii}}$ 和 $(x_{jk}-\bar{x}_k)/\sqrt{s_{kk}}$ 取代.标准化值是可通约的,因为两个集合都是中心在零,并用单位标准差表达.样本相关系数正是标准观测值的样本协方差.

尽管样本相关系数与样本协方差的记号相同,通常相关较容易解释,因为它的大小是有界的.概括地讲,样本相关系数 $r$ 有以下性质:

\begin{enumerate}[label=\arabic*.,leftmargin=*,itemsep=2pt]
\item $r$ 的值必定在 $-1$ 与 $+1$ 之间.
\item 这里 $r$ 度量的是线性结合的强度.如果 $r=0$,这就意味着分量之间无线性结合.另外,$r$ 的正负号指出了结合的方向; $r<0$ 意味着一个趋势,即当一对数中一个值大于它的平均值,另一个值则小于它的平均值; $r>0$ 表示这样的趋势,即当一对数中一个值大时另一个值也大,或者两值一起小.
\item 若第 $i$ 个变量的测量值变为 $y_{ji}=ax_{ji}+b, j=1,2,\cdots,n$,且第 $k$ 个变量的测量值变为 $y_{jk}=cx_{jk}+d, j=1,2,\cdots,n$,假定常数 $a$ 和 $c$ 的正负号相同,则 $r_{ik}$ 的值保持不变.
\end{enumerate}

参量 $s_{ik}$ 和 $r_{ik}$ 一般不能传送有关两个变量间结合的全部信息.可能存在不被描述统计量揭示的非线性结合.协方差和相关系数提供了线性结合或沿直线结合的度量.它们的值对其他类型的结合不能给出太多信息.另一方面,这些参量对于“杂乱的”观测值(“离群值”)非常敏感,可能会表示出事实上几乎不存在的结合.尽管有这些缺点,协方差和相关系数仍常常被



% UNCERTAIN: pg-013 章首页上方横排英文 "Applied Multivariate Statistical Analysis" 为章首装饰带（版式元素，非正文），未转写；如主文档需要，请指示。
% UNCERTAIN: pg-013 末行 "…分析类型的采" 与 pg-014 首行 "用很大程度上…" 跨页断词，已用行尾 % 接续为 "采用"，请核对。
% UNCERTAIN: pg-018 末行 "…协方差和相关系数仍常常被" 为页末断句，续文在 pg-019（超出本次范围）。
% UNCERTAIN: pg-016 "数据" 二字在原书为黑体（"(通常称为数据)"），已按 \mvbold{数据} 处理。
% UNCERTAIN: pg-017 例 1.1 矩阵右侧的实心方块 ■ 为原书"例结束"符号，已用 \blacksquare 还原。
% UNCERTAIN: pg-015 "cDNA 微序列实验" 一词已放大核对，字形为"微序列"（非"微阵列"），存疑待复核。
